\section{Separable evaluation images}
\label{sec:evaluation}

The positive and negative arguments will use the same test for the existence
of a separable quotient.  Let $X$ be a Banach space over
$\mathbb K\in\{\R,\C\}$, and write $B_X$ for its closed unit ball.
For a closed subspace $F\subset X^*$, define the evaluation operator
\[
 R_F:X\longrightarrow F^*,\qquad (R_Fx)(f)=f(x).
\]
The relevant condition is norm separability of $R_F(X)$, rather than
separability of $F$ itself.  We first explain why this condition yields an
actual quotient.

\begin{lemma}[Prefix normers and surjectivity]\label{ev:prefix}
Let $(f_n)_{n\ge1}$ be a sequence in $X^*$ with
$0<a\le\|f_n\|\le b<\infty$.  Suppose that, for each $m\ge1$, a finite
set $D_m\subset B_X$ satisfies
\begin{equation}\label{ev:normers}
 \max_{d\in D_m}|v(d)|\ge\tfrac12\|v\|
 \quad\bigl(v\in\operatorname{span}\{f_1,\ldots,f_m\}\bigr),
 \qquad f_n(D_m)=\{0\}\quad(n>m).
\end{equation}
Then $(f_n)$ is a basic sequence whose prefix projections have norm at most
$2$, and its coefficients satisfy
\begin{equation}\label{ev:coefficients}
 |c_n|\le\frac4a\left\|\sum_i c_i f_i\right\|
 \quad\text{for every finite sum.}
\end{equation}
Put $G=\overline{\operatorname{span}}\{f_n:n\ge1\}$, let $f_n^\#$ be its
basis coefficient functionals, and put
$Z=\overline{\operatorname{span}}\{f_n^\#:n\ge1\}\subset G^*$.
If $R_G(X)\subset Z$, then $R_G:X\to Z$ is onto.  More precisely, its
adjoint satisfies
\[
 \|R_G^*h\|\ge\tfrac12\|h\|\qquad(h\in Z^*).
\]
\end{lemma}

\begin{proof}
The vectors are linearly independent: a vector $f_n$ in the preceding
span would vanish on its norming set $D_{n-1}$ and hence be zero.
For $v=\sum_{i=1}^N c_i f_i$ and $m\le N$, the values of $v$ on $D_m$
equal those of its prefix.  Thus
\[
 \left\|\sum_{i=1}^m c_i f_i\right\|
 \le2\max_{d\in D_m}|v(d)|\le2\|v\|.
\]
The prefix projections extend to $G$ and converge to the identity there,
because they do so on the dense subspace of finite sums.  Taking differences
of two consecutive prefixes gives \eqref{ev:coefficients}.

Write $T=R_G:X\to Z$.  For $h\in Z^*$, set
\[
 v_m=\sum_{i=1}^m h(f_i^\#)f_i\in G.
\]
If $d\in D_m$, then
$Td=\sum_{i=1}^m f_i(d)f_i^\#$, so
$(T^*h)(d)=v_m(d)$.  Consequently
$\|T^*h\|\ge\|v_m\|/2$.
For every finite linear combination $u$ of the coefficient functionals,
we have $h(u)=u(v_m)$ for all sufficiently large $m$.  Since these
combinations are dense in $Z$,
\[
 \|h\|\le\sup_m\|v_m\|\le2\|T^*h\|.
\]
This argument does not require the vectors $v_m$ to converge.

We recall why an adjoint lower bound gives surjectivity.  If
$\|T^*z^*\|\ge c\|z^*\|$ for every $z^*\in Z^*$, Hahn--Banach separation
of the balanced convex set $\overline{T(B_X)}$ gives
$cB_Z\subset\overline{T(B_X)}$.
For a residual $z\ne0$, choose $x\in X$ with
$\|x\|\le\|z\|/c$ and $\|z-Tx\|\le\|z\|/2$.
Repeating this step makes the residuals decrease geometrically and the
chosen vectors form an absolutely convergent series in $X$.
Its sum maps to the original $z$ and has norm at most $2\|z\|/c$.
Apply this with $c=1/2$.
\end{proof}

\begin{proposition}[Evaluation criterion]\label{ev:criterion}
For an infinite-dimensional real or complex Banach space $X$, the following
are equivalent:
\begin{enumerate}
\item $X$ admits a bounded linear surjection onto a separable
      infinite-dimensional Banach space.
\item There is a closed infinite-dimensional subspace $F\subset X^*$
      for which $R_F(X)$ is norm separable.
\end{enumerate}
In the second assertion, $F$ may be required to be separable.
\end{proposition}

\begin{proof}
Suppose that $T:X\to Y$ is such a surjection.  Choose a closed separable
infinite-dimensional subspace $U\subset Y^*$.  The open mapping theorem
makes $T^*$ bounded below, so $F=T^*U$ is a closed separable
infinite-dimensional subspace of $X^*$.  Under the isomorphism
$T^*|_U:U\to F$, evaluation at $x$ becomes evaluation at $Tx$.
The evaluation map $Y\to U^*$ is bounded and $Y$ is separable; hence
$R_F(X)$ is norm separable.

Conversely, assume the second assertion and choose $x_j\in X$ so that
$\{R_Fx_j:j\ge1\}$ is norm dense in $R_F(X)$.
We construct $f_n\in F$ of norm one.  At stage $n$, require $f_n$ to
annihilate $x_1,\ldots,x_n$ and all finite sets $D_m$ already chosen,
$m<n$.  These conditions define a finite-codimensional subspace of $F$,
so a unit vector satisfying them exists.  After choosing $f_n$, choose
$D_n\subset B_X$ that is $1/2$-norming for
$\operatorname{span}\{f_1,\ldots,f_n\}$.  Such a finite set exists by
compactness of the finite-dimensional unit sphere and the definition of
the dual norm.

Use the spaces $G$ and $Z$ from Lemma~\ref{ev:prefix}.  For each $j$,
\[
 R_Gx_j=\sum_{i<j}f_i(x_j)f_i^\#\in Z.
\]
Restriction from $F^*$ to $G^*$ is contractive.  Density of the evaluations
at the $x_j$, and closedness of $Z$, therefore give $R_G(X)\subset Z$.
Lemma~\ref{ev:prefix} makes $R_G:X\to Z$ onto.  The coefficient
functionals are linearly independent, so $Z$ is separable and
infinite-dimensional.  The space $G$ also supplies the separable witness
in the second assertion.
\end{proof}

\subsection{Small restrictions of strictly singular operators}

We will use one elementary selection principle both in changing scalars
and in the later construction on arbitrary subspaces.  An operator is
\emph{strictly singular} if its restriction to no closed
infinite-dimensional subspace is bounded below.

\begin{lemma}\label{ev:small}
Let $T:Y\to Z$ be a strictly singular operator between real or complex
Banach spaces.  For every closed infinite-dimensional $M\subset Y$ and
every $\varepsilon>0$, there is a closed infinite-dimensional
$N\subset M$ such that $\|T|_N\|\le\varepsilon$.
For a finite family of strictly singular operators on $Y$, the same $N$
can be chosen to satisfy this bound for every operator in the family.
\end{lemma}

\begin{proof}
Inductively choose unit vectors $y_n\in M$ and finite $1/2$-norming sets
in $B_{Y^*}$ for their prefix spans.  Require each new vector to annihilate
all earlier norming sets and to satisfy
$\|Ty_n\|<\varepsilon 2^{-n}/4$.
The annihilation conditions leave an infinite-dimensional closed
subspace; strict singularity says that $T$ is not bounded below there,
so the choice is possible.  The prefix calculation in
Lemma~\ref{ev:prefix}, with vectors and norming functionals interchanged,
shows that $(y_n)$ is basic and
$|a_n|\le4\|\sum_i a_i y_i\|$.  Therefore
\[
 \left\|T\sum_n a_n y_n\right\|
 \le4\sum_n\|Ty_n\|\left\|\sum_n a_n y_n\right\|
 \le\varepsilon\left\|\sum_n a_n y_n\right\|.
\]
Continuity gives the bound on
$N=\overline{\operatorname{span}}\{y_n:n\ge1\}$.
For finitely many operators, apply the result successively inside the
subspace already obtained.  Restrictions of strictly singular operators
remain strictly singular, and previous bounds persist on smaller subspaces.
\end{proof}
